\documentclass[11pt]{article}
\usepackage[margin=1in]{geometry}
\usepackage{amsmath,amssymb,amsthm}
\usepackage{xurl}
\usepackage{hyperref}
\sloppy
\let\oldus\_ \renewcommand{\_}{\oldus\allowbreak}
\newtheorem{theorem}{Theorem}
\newtheorem{lemma}[theorem]{Lemma}
\theoremstyle{definition}
\newtheorem{definition}[theorem]{Definition}
\newcommand{\lean}[1]{\texttt{#1}}

\title{Conjecture 00000002511: symmetric tensors are homogeneous polynomials,\\ and the dictionary is multiplicative}
\author{Jackmeson1}
\date{}

\begin{document}
\maketitle

\section{Statement and reading}

\textbf{Conjecture (verbatim).} \emph{Definition: Symmetric tensors: the symmetric version.
Conjecture: The symmetric version of tensors corresponds to homogeneous polynomials: symmetric
tensors are forms, and the dictionary closes under homogeneous polynomials.}

\textbf{Reading.} This is the classical dictionary between symmetric tensors and homogeneous
polynomials (forms); we prove it as a known result, it is not new. We fix a field $K$ of
characteristic $0$, a $K$-vector space $V$ with a finite basis $b=(b_j)_{j\in\iota}$
($n=|\iota|$), and a degree $d\ge 0$. We prove:
\begin{enumerate}
\item[(a)] \emph{Symmetric tensors are forms:} the dictionary map sends every symmetric
  $d$-tensor to a homogeneous polynomial of degree $d$ in the $n$ variables $X_j$.
\item[(b)] \emph{Correspondence:} restricted to symmetric tensors, the dictionary is a linear
  bijection $\mathrm{Sym}^d(V)\to K[X_j: j\in\iota]_d$.
\item[(c)] \emph{The dictionary closes under products of forms:} the symmetric product
  $T\odot U$ of symmetric tensors is sent to the product of the two forms, and conversely the
  product $pq$ of two forms corresponds to the symmetric product of the tensors of $p$ and $q$.
\end{enumerate}
\textbf{Conventions.} (i) Characteristic $0$ is assumed (the standard setting; both sources
below work in it). We do not consider other characteristics. (ii) $d=0$ is allowed:
$V^{\otimes 0}=K$ and the forms of degree $0$ are the constants. (iii) Product convention: the
symmetrisation is the \emph{averaging} operator $\mathrm{Sym}\,t=\frac1{d!}\sum_{\sigma\in S_d}
\sigma\cdot t$ and $T\odot U=\mathrm{Sym}(T\otimes U)$, where $T\otimes U$ puts the $d$ factors of
$T$ before the $e$ factors of $U$. Only this normalisation is formalised. (iv) The polynomials are formal polynomials in the
coordinates $X_j$; Theorem~\ref{thm:eval} identifies them with the polynomial functions
$\varphi\mapsto T(\varphi,\dots,\varphi)$ on $V^*$.

\section{Definitions (as formalised)}

\begin{definition}
$V^{\otimes d}$ is Mathlib's \lean{PiTensorProduct} $\bigotimes_{k\in\mathrm{Fin}\,d}V$
(\lean{TPow K V d}). A permutation $\sigma\in S_d$ acts by permuting the tensor factors,
$\sigma\cdot(v_1\otimes\dots\otimes v_d)=v_{\sigma^{-1}(1)}\otimes\dots\otimes v_{\sigma^{-1}(d)}$
(\lean{permAct}, which is \lean{PiTensorProduct.reindex}). The \emph{symmetric tensors}
$\mathrm{Sym}^d(V)$ (\lean{symTensors K V d}) are the tensors fixed by every $\sigma\in S_d$.
\end{definition}
This is the definition in the retrieved Wikipedia article: ``Then T is a symmetric tensor if
[$\tau_\sigma T=T$] for the braiding maps associated to every permutation $\sigma$ on the symbols
\{1,2,...,k\}''.

\begin{definition}
The generic linear form is $\ell(v)=\sum_j b^*_j(v)\,X_j$, i.e.\ the linear map with
$b_j\mapsto X_j$ (\lean{linForm}). The dictionary map $P=P_{b,d}:V^{\otimes d}\to K[X_j:j\in\iota]$
(\lean{toPoly}) is the linear map with $P(v_1\otimes\dots\otimes v_d)=\ell(v_1)\cdots\ell(v_d)$
(defined by the universal property \lean{PiTensorProduct.lift}). Forms of degree $d$ are Mathlib's
\lean{MvPolynomial.homogeneousSubmodule $\iota$ K d}.
\end{definition}

\begin{theorem}[\lean{eval\_toPoly}]\label{thm:eval}
For every linear functional $\varphi\in V^*$ and $T\in V^{\otimes d}$,
$P(T)\big((\varphi(b_j))_j\big)=(\varphi\otimes\dots\otimes\varphi)(T)$, i.e.\ $P(T)$ is the
polynomial $\varphi\mapsto T(\varphi,\dots,\varphi)$ written in the coordinates of $\varphi$.
\end{theorem}
\begin{proof}
Both sides are linear in $T$; on $v_1\otimes\dots\otimes v_d$ both equal $\prod_k\varphi(v_k)$,
since evaluation at $(\varphi(b_j))_j$ composed with $\ell$ is $\varphi$ (check on the basis).
\end{proof}

\section{Proof}

Let $B_i=b_{i_1}\otimes\dots\otimes b_{i_d}$, $i:\mathrm{Fin}\,d\to\iota$, be the product basis of
$V^{\otimes d}$ (\lean{tBasis}, Mathlib's \lean{Basis.piTensorProduct}) and $T_i$ the coordinates of
$T$. The \emph{content} of a word $i$ is the exponent vector $c(i)=\sum_k e_{i_k}\in\mathbb N^\iota$
(\lean{content}); $c(i)_a=\#\{k: i_k=a\}$ and $|c(i)|=d$.

\begin{lemma}\label{lem:coord}
(1) $P(B_i)=X^{c(i)}$, hence $P(T)=\sum_i T_i X^{c(i)}$ and the coefficient of $X^\alpha$ in $P(T)$
is $\sum_{c(i)=\alpha}T_i$ (\lean{toPoly\_tBasis}, \lean{toPoly\_eq\_sum}, \lean{coeff\_toPoly}).
(2) $(\sigma\cdot T)_i=T_{i\circ\sigma}$ (\lean{repr\_permAct}).
(3) If $c(i)=c(j)$ then $i=j\circ\sigma$ for some $\sigma\in S_d$ (\lean{exists\_perm\_of\_content\_eq}:
the fibres $\{k:i_k=a\}$ and $\{k:j_k=a\}$ have equal size, and \lean{Equiv.ofFiberEquiv} glues
bijections between them). Hence a symmetric $T$ has $T_i=T_j$ whenever $c(i)=c(j)$
(\lean{repr\_sym\_eq}).
(4) Every $\alpha\in\mathbb N^\iota$ with $|\alpha|=d$ is $c(i)$ for some $i$
(\lean{exists\_content\_eq}, induction on $d$, removing one unit from a nonzero entry).
\end{lemma}

\emph{(a) Forms.} By Lemma~\ref{lem:coord}(1), $P(T)$ is a combination of monomials of degree
$d$, so $P(T)$ is homogeneous of degree $d$ for every $T$, in particular for symmetric $T$
(\lean{toPoly\_mem}).

\emph{(b) Injectivity.} Let $T$ be symmetric with $P(T)=0$ and fix $i$. With $F=\{j:c(j)=c(i)\}$,
Lemma~\ref{lem:coord}(1),(3) give $0=\sum_{j\in F}T_j=|F|\,T_i$. Since $i\in F$, $|F|\ge1$, and
$|F|\neq0$ in $K$ because $\mathrm{char}\,K=0$; so $T_i=0$ (\lean{toPoly\_eq\_zero}).

\emph{(b) Surjectivity.} Let $|\alpha|=d$, $F_\alpha=\{j:c(j)=\alpha\}$ (nonempty by
Lemma~\ref{lem:coord}(4)) and $S_\alpha=|F_\alpha|^{-1}\sum_{j\in F_\alpha}B_j$. By
Lemma~\ref{lem:coord}(2) and $c(j\circ\sigma)=c(j)$, $S_\alpha$ is symmetric, and by (1)
$P(S_\alpha)=X^\alpha$ (\lean{exists\_sym\_monomial}). The forms of degree $d$ are spanned by the
monomials $X^\alpha$ with $|\alpha|=d$ (Mathlib \lean{homogeneousSubmodule\_eq\_finsupp\_supported}),
so every form of degree $d$ is $P$ of a symmetric tensor (\lean{exists\_sym\_of\_mem}). Hence the
restriction \lean{symToForm b d} is bijective (\lean{symToForm\_bijective}) and gives the linear
isomorphism \lean{symEquiv b d} $:\mathrm{Sym}^d(V)\simeq K[X]_d$; in particular the two spaces have
equal dimension (\lean{finrank\_symTensors}).

\emph{(c) Products.} Let $T\otimes U\in V^{\otimes(d+e)}$ be
\lean{reindex finSumFinEquiv (tmulEquiv (T $\otimes$ U))} (\lean{tmulL}). On pure tensors
$P(T\otimes U)=\prod_{k<d+e}\ell(\cdot)=P(T)P(U)$ (split the product with \lean{Fin.prod\_univ\_add}),
so by bilinearity $P(T\otimes U)=P(T)P(U)$ (\lean{toPoly\_tmulL}). Also $P(\sigma\cdot t)=P(t)$
(the product of the $\ell(v_k)$ is reordered, \lean{toPoly\_permAct}), hence $P(\mathrm{Sym}\,t)=
\frac1{d!}\cdot d!\cdot P(t)=P(t)$ (\lean{toPoly\_symz}). We check that $\mathrm{Sym}\,t$ is symmetric
(\lean{symz\_mem}) and that $\mathrm{Sym}$ fixes symmetric tensors (\lean{symz\_of\_mem}). Therefore
$P(T\odot U)=P(T\otimes U)=P(T)P(U)$ (\lean{symToForm\_symProd}). Applying the inverse isomorphisms
gives the converse statement in (c).

\section{Lean formalisation}

File \lean{Conjecture2511/Basic.lean} (namespace \lean{C2511}, 341 lines, only
\lean{import Mathlib}). The main theorem is
\begin{verbatim}
theorem symmetric_tensors_are_forms [CharZero K] (b : Module.Basis iota K V) :
  (forall (d : Nat) (T : TPow K V d), T mem symTensors K V d ->
     toPoly b d T mem homogeneousSubmodule iota K d) /\
  (forall d : Nat, Function.Bijective (symToForm b d)) /\
  (forall (d e : Nat) (T : symTensors K V d) (U : symTensors K V e),
     (symToForm b (d + e) (symProd T U) : MvPolynomial iota K) =
       symToForm b d T * symToForm b e U) /\
  (forall (d e : Nat) (p : homogeneousSubmodule iota K d)
     (q : homogeneousSubmodule iota K e),
     (symEquiv b (d + e)).symm <p * q, p.2.mul q.2> =
       symProd ((symEquiv b d).symm p) ((symEquiv b e).symm q))
\end{verbatim}
(ASCII transliteration of the Lean source; $\iota$ is an arbitrary \lean{Fintype}, $V$ any $K$-vector space with basis $b$). Supporting
results: \lean{eval\_toPoly}, \lean{symEquiv}, \lean{finrank\_symTensors}, \lean{symz\_of\_mem}.
No component from earlier packages was reused.

\textbf{Axiom audit.} \lean{lake build} succeeds; \lean{\#print axioms} for
\lean{symmetric\_tensors\_are\_forms}, \lean{eval\_toPoly}, \lean{symEquiv},
\lean{finrank\_symTensors} and \lean{symz\_of\_mem} each reports
\lean{[propext, Classical.choice, Quot.sound]}. There is no \lean{sorry}, no \lean{native\_decide}
and no added axiom.

\section{Sources (retrieved)}
\begin{itemize}
\item Wikipedia, \emph{Symmetric tensor}, \url{https://en.wikipedia.org/wiki/Symmetric_tensor}:
``The space of symmetric tensors of order r on a finite-dimensional vector space V is naturally
isomorphic to the dual of the space of homogeneous polynomials of degree r on V.  Over fields of
characteristic zero, the graded vector space of all symmetric tensors can be naturally identified
with the symmetric algebra on V.'' The same article defines the averaging symmetrisation and the
product: ``Given two tensors T1 $\in$ Symk1(V) and T2 $\in$ Symk2(V), we use the symmetrization
operator to define:'' $T_1\odot T_2=\mathrm{Sym}(T_1\otimes T_2)$.
\item P. Comon, G. Golub, L.-H. Lim, B. Mourrain, \emph{Symmetric tensors and symmetric tensor
rank}, SIAM J. Matrix Anal. Appl. (2008), \url{https://arxiv.org/abs/0802.1681}, Section 3.1
``Equivalence with homogeneous polynomials'': ``every symmetric tensor of order and dimension may
be uniquely associated with a homogeneous polynomial of degree in variables'' and ``the
correspondence between symmetric tensors and homogeneous polynomials is obviously bijective''
(mathematical symbols were lost in the retrieved text).
\end{itemize}

\end{document}
